\documentclass[12pt]{amsart}

\usepackage{amsmath,amssymb,amsthm}
\usepackage{dsfont}
\usepackage[colorlinks=false,pdfborder={0 0 0}]{hyperref}
\usepackage{booktabs}

% ── Theorem environments ──
\newtheorem{theorem}{Theorem}
\newtheorem{lemma}[theorem]{Lemma}
\newtheorem{proposition}[theorem]{Proposition}
\newtheorem{corollary}[theorem]{Corollary}

\theoremstyle{definition}
\newtheorem*{axiom}{Axiom}

\theoremstyle{remark}
\newtheorem*{remark}{Remark}

% ── Operators ──
\DeclareMathOperator{\supp}{supp}
\DeclareMathOperator{\Var}{Var}

\subjclass[2020]{60G57, 60E10, 28A80, 76F55}
\keywords{Multiplicative cascades, log-Poisson distribution,
hierarchical symmetry, multifractal spectrum, Hausdorff moment problem}

\begin{document}

\title[Hierarchical symmetry selects log-Poisson cascades]%
{Hierarchical symmetry selects log-Poisson cascades: \\
classification, uniqueness, and stability}

\author{E.~M.~Freeburg}
\address{Independent Researcher}
\date{}

\begin{abstract}
Within i.i.d.\ multiplicative cascades, a single axiom---the
hierarchical symmetry, a linear contraction on incremental scaling
exponents---is shown to be necessary and sufficient for the cascade
multiplier to be log-Poisson.  We establish three results:
(1)~a characterization theorem proving that the hierarchical symmetry
uniquely determines the log-Poisson distribution with explicit
parameters; (2)~a classification theorem proving that the hierarchical
symmetry selects exactly the log-Poisson class from the full
log-infinitely-divisible family, excluding log-normal, log-stable, and
all intermediate generators; and (3)~a stability theorem proving that
approximate hierarchical symmetry implies approximate log-Poisson,
with an explicit $O(\sqrt{\varepsilon})$ Wasserstein bound.  The
proofs reduce the problem to the Hausdorff moment problem on~$[0,1]$
via the change of variables $u = e^{kx}$, where determinacy and
stability follow from classical results.
\end{abstract}

\maketitle

% ══════════════════════════════════════════════════════════════════════
\section{Introduction}
% ══════════════════════════════════════════════════════════════════════

Multiplicative cascades model the successive fragmentation of a
conserved quantity across scales and arise in fully developed
turbulence \cite{K62,Man74}, rainfall, finance, and other settings
exhibiting intermittent, scale-invariant fluctuations.  The
mathematical foundations of multiplicative cascades were established
by Kahane and Peyri\`ere~\cite{KP76}; see also Barral and
Mandelbrot~\cite{BM02} for later developments.  The statistical
properties of a cascade are encoded in the scaling exponents
$\zeta_p$ of the structure functions $S_p(\ell) = \langle
|\Phi(\ell)|^p \rangle \sim \ell^{\zeta_p}$.  A central question is:
\emph{which probability distributions on the cascade multiplier $W$
are compatible with observed scaling laws?}

Kolmogorov~\cite{K62} proposed log-normal multipliers, leading to
quadratic scaling exponents.  Z.-S.~She and
L\'ev\^eque~\cite{SL94} introduced a hierarchical symmetry for the
scaling exponents and derived a different, log-Poisson exponent
formula that has since shown excellent agreement with experimental
data.  Dubrulle~\cite{Dub94} independently identified the log-Poisson
form.  Z.-S.~She and Waymire~\cite{SW95} used the
L\'evy--Khintchine representation to argue that this symmetry selects
log-Poisson within the log-infinitely-divisible family; Dubrulle and
Graner~\cite{DG96} reached a similar conclusion via symmetry groups.
These works provided compelling physical arguments but did not supply
rigorous proofs.  Z.-S.~She and Zhang~\cite{SZ09} subsequently
proposed that the hierarchical symmetry is \emph{universal}---applicable
not only to turbulence but to general multi-scale fluctuation systems
including MHD turbulence, natural image statistics, and biological
signals---and should serve as a standard analytical framework.  The
present paper supplies the rigorous mathematical foundation for this
program.

The present paper provides the rigorous mathematical treatment.  We
formalize the hierarchical symmetry as a single axiom~(A1) and prove
three main results.

\emph{Characterization}
(Theorem~\ref{thm:characterization}).  A1 uniquely determines the
cascade multiplier $W$ to be log-Poisson, with parameters
$(\,a = \gamma\ln r,\; b = (\ln\beta)/k,\; \lambda = -C\ln r\,)$
expressed in terms of the observable scaling exponents.  No other
distribution is compatible with~A1.

\emph{Classification}
(Theorem~\ref{thm:classification}).  Within the full
log-infinitely-divisible family, A1 selects exactly the log-Poisson
class.  No log-normal, log-stable, or intermediate generator
satisfies~A1.

\emph{Stability}
(Theorem~\ref{thm:stability}).  If A1 holds only approximately, with
residuals bounded by~$\varepsilon$, then the cascade multiplier
distribution is within $O(\sqrt{\varepsilon})$ of log-Poisson in the
Wasserstein-1 metric, with an explicit, computable constant.

The converse---that log-Poisson multipliers imply
A1---is established in Proposition~\ref{prop:converse}, yielding a
biconditional equivalence (Corollary~\ref{cor:biconditional}).

\smallskip
\noindent\textbf{Relation to prior work.}
The exponent formula~\eqref{eq:zeta} (Lemma~\ref{lem:exponent},
steps~(1)--(3)) and the log-Poisson identification are due to
Z.-S.~She and L\'ev\^eque~\cite{SL94} and Dubrulle~\cite{Dub94}.
Z.-S.~She and Waymire~\cite{SW95} gave the first argument connecting
A1 to the L\'evy--Khintchine classification.  The following results are new: the moment-determinacy proof via Carleman's
condition (Lemma~\ref{lem:determinacy}); the converse
(Proposition~\ref{prop:converse}) and biconditional equivalence
(Corollary~\ref{cor:biconditional}); the full Log-ID Classification
Theorem via the Hausdorff moment problem
(Theorem~\ref{thm:classification}); the determinacy dichotomy
(Proposition~\ref{prop:determinacy}); and the stability theorem with
explicit $O(\sqrt{\varepsilon})$ bound
(Theorem~\ref{thm:stability}).

\smallskip
\noindent\textbf{Method.}
The key technique is the change of variables $u = e^{kx}$, which maps
the L\'evy measure from $(-\infty,0]$ to the compact interval
$[0,1]$.  On $[0,1]$ the Hausdorff moment problem is automatically
determinate and stable, and the entire classification and stability
theory follows from the Weierstrass approximation theorem,
Chebyshev's inequality, and explicit coupling constructions.

\smallskip
Throughout this paper, $\log$ denotes the natural logarithm.


% ══════════════════════════════════════════════════════════════════════
\section{Setup}
\label{sec:setup}
% ══════════════════════════════════════════════════════════════════════

Let $r \in (0,1)$ be a scale ratio.  A \emph{multiplicative cascade}
generates a random positive measure $\mu$ on nested sets $B_0 \supset
B_1 \supset \cdots$ via
\[
  \mu(B_{n+1}) = W_{n+1} \cdot \mu(B_n),
\]
where $\{W_n\}$ are i.i.d.\ positive random variables with
$\mathbb{E}[W] = 1$ (conservation of mean).

Let $\Phi(\ell)$ be the cascade observable at scale $\ell = r^n$.
Define the \emph{structure functions}
\[
  S_p(\ell) = \langle |\Phi(\ell)|^p \rangle = \ell^{\zeta_p},
\]
where $\zeta_p$ are the \emph{scaling exponents}, with $\zeta_0 = 0$.
The relation $\mathbb{E}[W^p] = r^{\zeta_p}$ identifies the
single-step moment structure of the multiplier with the observable's
scaling.

For a fixed integer $k \geq 1$ (the \emph{hierarchy step}), define
the \emph{moment ratios}
\[
  H_p(\ell) = \frac{S_{p+k}(\ell)}{S_p(\ell)} = \ell^{\delta_p},
\]
where $\delta_p = \zeta_{p+k} - \zeta_p$ are the \emph{incremental
exponents} at step~$k$.


% ══════════════════════════════════════════════════════════════════════
\section{The axiom}
\label{sec:axiom}
% ══════════════════════════════════════════════════════════════════════

\begin{axiom}[A1: Hierarchical Symmetry]
There exists $\beta \in (0,1)$ such that for all $p \in k\mathbb{N}_0$
(non-negative integer multiples of~$k$), the incremental exponents
satisfy
\begin{align}\label{eq:A1}
  \delta_{p+k}
  = (1 - \beta)\,\delta_\infty + \beta\,\delta_p.
  \tag{$*$}
\end{align}
where $\delta_\infty = \lim_{p \to \infty} \delta_p$ exists and is
finite.
\end{axiom}

A1 determines the full parameter set from the observable exponents:
\begin{center}\small
\begin{tabular}{@{}lll@{}}
\toprule
Parameter & Determined by & Meaning \\
\midrule
$\beta$ & Contraction ratio of~\eqref{eq:A1} & Coupling strength \\
$\gamma$ & $\delta_\infty / k$ & Linear drift \\
$C$ & $(\delta_0 - \delta_\infty)/(1 - \beta)$ & Concentration amplitude \\
\bottomrule
\end{tabular}
\end{center}

\smallskip\noindent\textit{Edge case} ($C = 0$).  If $\delta_0 =
\delta_\infty$, then $C = 0$ and $\zeta_p = \gamma p$ (monofractal
scaling).  The cascade multiplier $W = r^\gamma$ is deterministic.
A1 is trivially satisfied for any $\beta \in (0,1)$.  The
classification and stability theorems assume $C > 0$ (nontrivial
intermittency).


% ══════════════════════════════════════════════════════════════════════
\section{Results}
\label{sec:results}
% ══════════════════════════════════════════════════════════════════════


% ── Lemma 1 ─────────────────────────────────────────────────────────
\begin{lemma}[Exponent Form]\label{lem:exponent}
If the incremental exponents $\{\delta_p\}$ satisfy the A1
recurrence~\eqref{eq:A1} with $\beta \in (0,1)$, then with $\zeta_0
= 0$:
\begin{align}\label{eq:zeta}
  \zeta_p = \gamma p + C\bigl(1 - \beta^{p/k}\bigr),
  \tag{$**$}
\end{align}
where $\gamma = \delta_\infty / k$ and $C = (\delta_0 -
\delta_\infty)/(1 - \beta)$.

We emphasize that this lemma is purely algebraic and involves no
probabilistic content.
\end{lemma}

\begin{proof}
(1)~The recurrence~\eqref{eq:A1} is first-order linear with fixed
point $\delta_\infty$:
\[
  \delta_{p+k} - \delta_\infty = \beta\,(\delta_p - \delta_\infty).
\]

(2)~At $p = mk$ (integer multiples of~$k$), iteration gives
\[
  \delta_{mk} - \delta_\infty
  = (\delta_0 - \delta_\infty)\,\beta^m.
\]
Replacing $m = p/k$:
\[
  \delta_p
  = \delta_\infty + (\delta_0 - \delta_\infty)\,\beta^{p/k}.
\]
(This formula is derived at $p \in k\mathbb{N}_0$.  For general $p
\geq 0$, we define $\zeta_p$ by~\eqref{eq:zeta}; the function
$\beta^{p/k}$ is well-defined for all real $p \geq 0$ since $\beta >
0$.  The moment-based arguments in Lemma~\ref{lem:determinacy} and
Theorem~\ref{thm:characterization} require only integer~$p$, where
the formula is proved.)

(3)~With $\zeta_0 = 0$, sum the step-$k$ increments using the
geometric series:
\[
  \zeta_p
  = \frac{p}{k}\,\delta_\infty
    + \frac{\delta_0 - \delta_\infty}{1 - \beta}\,
      \bigl(1 - \beta^{p/k}\bigr).
\]
Identifying $\gamma = \delta_\infty/k$ and $C = (\delta_0 -
\delta_\infty)/(1-\beta)$:
\[
  \zeta_p = \gamma p + C\bigl(1 - \beta^{p/k}\bigr).  \qedhere
\]
\end{proof}


% ── Lemma 2 ─────────────────────────────────────────────────────────
\begin{lemma}[Moment Determinacy]\label{lem:determinacy}
Let $\{W_n\}$ be an i.i.d.\ multiplicative cascade with scaling
exponents $\zeta_p = \gamma p + C(1 - \beta^{p/k})$ as in
Lemma~\ref{lem:exponent}.  Then the distribution of~$W$ is uniquely
determined by its moments.
\end{lemma}

\begin{proof}
(1)~By independence across cascade levels, $\mathbb{E}[(W_1 \cdots
W_n)^p] = (\mathbb{E}[W^p])^n$.  Combined with $S_p(\ell) =
\ell^{\zeta_p} = r^{n\zeta_p}$, this gives for a single step:
\[
  \mathbb{E}[W^p] = r^{\zeta_p}
  = e^{(\gamma \ln r)\,p}\cdot e^{C \ln r\,(1 - \beta^{p/k})}.
\]
The moments $\mathbb{E}[W^p]$ are finite and positive for all $p \geq
0$, since $r \in (0,1)$ and $\zeta_p$ is finite by
Lemma~\ref{lem:exponent}.

(2)~The distribution of~$W$ is uniquely determined by its moments if
\[
  \sum_{p=1}^{\infty}
  \bigl(\mathbb{E}[W^{2p}]\bigr)^{-1/(2p)} = \infty
  \qquad\text{(Carleman's condition).}
\]
Now
\[
  \bigl(\mathbb{E}[W^{2p}]\bigr)^{1/(2p)}
  = r^{\gamma + C(1 - \beta^{2p/k})/(2p)}
  \;\longrightarrow\; r^\gamma
  \quad\text{as } p \to \infty.
\]
Since $r \in (0,1)$ and $\gamma$ is finite, $r^{-\gamma} > 0$, so
the terms $(\mathbb{E}[W^{2p}])^{-1/(2p)}$ converge to the positive
constant $r^{-\gamma}$.  A series whose terms do not converge to zero
diverges; therefore the Carleman sum diverges.  The moments uniquely
determine~$W$.

(Since $W > 0$, this is the Stieltjes moment problem on $[0,\infty)$.
The Hamburger condition verified here implies the Stieltjes result
\textit{a fortiori}; see Akhiezer~\cite{Akh65}, Ch.~2, Theorem~2.1,
or Shohat and Tamarkin~\cite{ST43}, Ch.~II.)
\end{proof}


% ── Theorem 1 ───────────────────────────────────────────────────────
\begin{theorem}[Characterization]\label{thm:characterization}
Let $\{W_n\}$ be an i.i.d.\ multiplicative cascade whose incremental
scaling exponents satisfy A1.  Then:

\textup{(i)} \textbf{Scaling exponents.}
\[
  \zeta_p = \gamma p + C\bigl(1 - \beta^{p/k}\bigr),
\]
where $\gamma = \delta_\infty/k$ and $C = (\delta_0 -
\delta_\infty)/(1-\beta)$.

\textup{(ii)} \textbf{Uniqueness.}  The cascade multiplier $W$ is
uniquely determined to be log-Poisson:
\[
  \log W = a + bN, \qquad N \sim \mathrm{Poisson}(\lambda),
\]
\[
  a = \gamma \ln r, \quad
  b = \frac{\ln\beta}{k}, \quad
  \lambda = -C\ln r.
\]
No other probability distribution on~$W$ is compatible with A1.

\textup{(iii)} \textbf{Multifractal spectrum.}  Let $d$ denote the
spatial dimension of the cascade support.  Then
\[
  f(h) = d - C + Cx(1 - \ln x),
  \qquad
  x = \frac{k(h - \gamma)}{C|\ln\beta|},
\]
defined for $h \in [\gamma,\; \gamma + (C/k)|\ln\beta|]$.
\end{theorem}

\begin{proof}
\textit{(i)} Immediate from Lemma~\ref{lem:exponent}.

\textit{(ii)} By Lemma~\ref{lem:determinacy}, $W$ is uniquely
determined by its moments.  It remains to exhibit a distribution that
produces exactly these moments.  For $\log W = a + bN$ with $N \sim
\mathrm{Poisson}(\lambda)$:
\[
  \mathbb{E}[W^p] = e^{ap} \cdot \exp\bigl[\lambda(e^{bp} - 1)\bigr].
\]
Set $a = \gamma\ln r$, $b = (\ln\beta)/k$ so that $e^{bp} =
\beta^{p/k}$.  Then matching requires
\[
  \lambda(\beta^{p/k} - 1) = C|\ln r|\,(\beta^{p/k} - 1),
\]
giving $\lambda = -C\ln r > 0$.  This holds for all real $p \geq 0$
simultaneously, since the Poisson MGF is defined for all $t \in
\mathbb{R}$.

The log-Poisson with $(a,b,\lambda)$ above produces exactly the
moments from Lemma~\ref{lem:determinacy}.  By uniqueness, $W$ is
log-Poisson.  No other distribution is possible.

\textit{(iii)}  The singularity spectrum $f(h) = \inf_p\,[ph -
\zeta_p + d]$.  Setting the derivative to zero:
\[
  h - \gamma + \frac{C}{k}(\ln\beta)\,\beta^{p/k} = 0
  \quad\Longrightarrow\quad
  h - \gamma = \frac{C|\ln\beta|}{k}\,\beta^{p/k}.
\]
Define $x = \beta^{p/k} = k(h-\gamma)/(C|\ln\beta|)$.  Then $p =
-k\ln x / |\ln\beta|$ and
\[
  p(h - \gamma) = \frac{-k\ln x}{|\ln\beta|}
    \cdot \frac{C|\ln\beta|}{k}\,x = -Cx\ln x.
\]
Therefore
\[
  f(h) = d - C + Cx(1 - \ln x), \qquad
  x = \frac{k(h-\gamma)}{C|\ln\beta|}.
\]
Boundary checks: $p = 0 \Rightarrow x = 1 \Rightarrow f = d$; $p \to
\infty \Rightarrow x \to 0 \Rightarrow f \to d - C$.  Concavity:
$\zeta_p'' = -(C/k^2)(\ln\beta)^2\beta^{p/k} < 0$.
\end{proof}


\begin{remark}[Conservation]
The setup assumes $\mathbb{E}[W] = 1$, which requires $\zeta_1 = 0$.
Substituting into~\eqref{eq:zeta}: $\gamma + C(1 - \beta^{1/k}) =
0$, giving $\gamma = -C(1-\beta^{1/k})$.  This is a constraint
relating $\gamma$ to~$C$ and~$\beta$, reducing the free parameters
from three to two.
\end{remark}


% ── Proposition 1 ───────────────────────────────────────────────────
\begin{proposition}[Converse]\label{prop:converse}
If the cascade multiplier $W$ is log-Poisson---that is, $\log W = a +
bN$ with $N \sim \mathrm{Poisson}(\lambda)$, $b < 0$, $\lambda >
0$---then the incremental scaling exponents satisfy A1 with $\beta =
e^{bk} \in (0,1)$.
\end{proposition}

\begin{proof}
The moment generating function gives $\mathbb{E}[W^p] = \exp(ap +
\lambda(e^{bp} - 1))$, so $\zeta_p = \bigl(ap + \lambda(e^{bp} -
1)\bigr)/\ln r$.  The step-$k$ increments are
\[
  \delta_p = \frac{ak + \lambda e^{bp}(e^{bk}-1)}{\ln r}.
\]
Setting $\beta = e^{bk} \in (0,1)$ (since $b < 0$, $k \geq 1$):
\[
  \delta_p = \frac{ak}{\ln r}
  + \frac{\lambda(\beta - 1)}{\ln r}\,\beta^{p/k}.
\]
As $p \to \infty$: $\beta^{p/k} \to 0$, so $\delta_\infty = ak/\ln
r$.  The deviation is
\[
  \delta_p - \delta_\infty
  = \frac{\lambda(\beta - 1)}{\ln r}\,\beta^{p/k}.
\]
At $p + k$:
\[
  \delta_{p+k} - \delta_\infty
  = \frac{\lambda(\beta-1)}{\ln r}\,\beta^{(p+k)/k}
  = \beta\,(\delta_p - \delta_\infty).
\]
Therefore $\delta_{p+k} = (1-\beta)\delta_\infty + \beta\delta_p$,
which is exactly A1.
\end{proof}

\begin{corollary}[Biconditional]\label{cor:biconditional}
Within i.i.d.\ multiplicative cascades, A1 is necessary and
sufficient for log-Poisson:
\[
  \text{A1 holds}
  \;\;\Longleftrightarrow\;\;
  W \text{ is log-Poisson (with } b < 0\text{).}
\]
The forward direction is
Theorem~\textup{\ref{thm:characterization}(ii)}; the reverse is
Proposition~\textup{\ref{prop:converse}}.
\end{corollary}


% ── Theorem 2 ───────────────────────────────────────────────────────
\begin{theorem}[Log-ID Classification]\label{thm:classification}
Let $\{W_n\}$ be an i.i.d.\ multiplicative cascade with nontrivial
intermittency ($C > 0$), whose generator $\log W$ is infinitely
divisible with L\'evy triplet $(a, \sigma^2, \nu)$.  Then A1 holds
with $\beta \in (0,1)$ if and only if $\sigma^2 = 0$ and $\nu =
\lambda\delta_b$ for some $b < 0$, $\lambda > 0$.  That is:
\begin{center}
\emph{A1 selects exactly the log-Poisson class from the full
log-infinitely-divisible family.}
\end{center}
No other log-ID cascade---log-normal, log-stable, or any
intermediate---satisfies A1.
\end{theorem}

\begin{proof}
\textit{Reverse direction.} If $\nu = \lambda\delta_b$ with $b < 0$
and $\sigma^2 = 0$, then $\log W = a + bN$ with $N \sim
\mathrm{Poisson}(\lambda)$, and A1 holds by
Proposition~\ref{prop:converse}.

\textit{Forward direction.}  Assume A1 holds.  We show $\sigma^2 = 0$
and $\nu = \lambda\delta_b$.

\textit{Step~1.}  The cumulant generating function of $\log W$ is
\[
  \psi(p) = ap + \frac{\sigma^2 p^2}{2}
  + \int\bigl(e^{px} - 1 - px\,\mathds{1}_{|x|\leq 1}\bigr)
  \,\nu(dx).
\]
With $\zeta_p = \psi(p)/\ln r$ and $\delta_p = (\psi(p+k) -
\psi(p))/\ln r$, define $\phi(p) = \psi(p+k) - \psi(p)$:
\[
  \phi(p) = ak + \sigma^2 k\!\left(p + \tfrac{k}{2}\right)
  + \int e^{px}(e^{kx}-1)\,\nu(dx)
  - k\!\int x\,\mathds{1}_{|x|\leq 1}\,\nu(dx).
\]
A1 requires $\delta_\infty = \lim_{p\to\infty}\delta_p$ to exist and
be finite, i.e., $\lim_{p\to\infty}\phi(p)$ is finite.

\textit{Step~2} ($\sigma^2 = 0$).  The term $\sigma^2 kp$ in
$\phi(p)$ diverges as $p \to \infty$ whenever $\sigma^2 > 0$.  Since
$\phi$ must have a finite limit, $\sigma^2 = 0$.  \textit{This
eliminates all log-normal and mixed Gaussian-jump generators.}

\textit{Step~3} ($\supp(\nu) \subseteq (-\infty,0]$).  For $x > 0$:
$e^{kx} - 1 > 0$ and $e^{px} \to \infty$ as $p \to \infty$.  If
$\nu$ has any mass on $(0,\infty)$, then $\int_{(0,\infty)}
e^{px}(e^{kx}-1)\,\nu(dx) \to \infty$ by monotone convergence,
making $\phi(p)$ unbounded.  Therefore $\supp(\nu) \subseteq
(-\infty,0]$.  \textit{This eliminates all generators with positive
jumps.}

\textit{Step~4} ($\nu$ is a single Dirac mass).  With $\sigma^2 = 0$
and $\supp(\nu) \subseteq (-\infty,0]$, write $c_0 = ak - k\int
x\,\mathds{1}_{|x|\leq 1}\,\nu(dx)$ and
\[
  \phi(p) = c_0 + \int_{(-\infty,0)} e^{px}(e^{kx}-1)\,\nu(dx).
\]
For $x < 0$ and $p \geq 0$: $|e^{px}(e^{kx}-1)| \leq |e^{kx}-1| = 1
- e^{kx}$ (since $|e^{px}| \leq 1$).  This bound is
$\nu$-integrable: near $x = 0$, $1 - e^{kx} \sim k|x|$ and $\int
k|x|\,\nu(dx) < \infty$ follows from the finiteness of $\delta_0$;
for $|x| > 1$, the L\'evy condition gives $\nu((-\infty,-1)) <
\infty$ and the integrand is bounded by~1.  By dominated convergence,
the integral $\to 0$ as $p \to \infty$, so $\phi_\infty = c_0$.

A1 at $p = mk$ gives $\phi(mk) - \phi_\infty = A\beta^m$, where $A =
(\delta_0 - \delta_\infty)\ln r$:
\begin{equation}\label{eq:dagger}
  \int_{(-\infty,0)} e^{mkx}(e^{kx}-1)\,\nu(dx) = A\beta^m
  \qquad\text{for all } m \geq 0.
\end{equation}
Substitute $u = e^{kx}$, mapping $(-\infty,0) \to (0,1)$.  Let
$\tilde\nu$ be the pushforward of $\nu$ under $x \mapsto e^{kx}$, and
define the signed measure $d\rho = (u-1)\,d\tilde\nu$ on $(0,1)$.
Condition~\eqref{eq:dagger} becomes
\[
  \int_{(0,1)} u^m\,d\rho(u) = A\beta^m
  \qquad\text{for all } m \geq 0.
\]
The $m = 0$ case gives $\int(u-1)\,d\tilde\nu = A$, so $\int(1-u)\,
d\tilde\nu = -A = |A| < \infty$.  Since $1 - u > 0$ on $(0,1)$, this
shows $\rho$ has finite total variation $|\rho|((0,1)) = |A|$, hence
is a finite signed measure on~$[0,1]$.

(The measure $\rho$ is defined on the open interval $(0,1)$.  It
extends to $[0,1]$ with $\rho(\{0\}) = \rho(\{1\}) = 0$: $\nu$ has no
atom at $x = 0$ by the L\'evy measure convention, so $\tilde\nu$ has
no atom at $u = 1$; and $x = -\infty$ maps to $u = 0$, which is not a
point of the measure.)

The measure $A\delta_\beta$ on $(0,1)$ has the same moments: $\int
u^m\,A\,d\delta_\beta = A\beta^m$.  Since polynomials are uniformly
dense in $C([0,1])$ (Weierstrass approximation theorem), moments
uniquely determine finite signed measures on $[0,1]$ (via the Riesz
representation theorem).  Therefore $\rho = A\delta_\beta$.

Recovering $\nu$: at $u = \beta$ (i.e., $x = b = (\ln\beta)/k$), we
have $(\beta - 1)\tilde\nu(\{\beta\}) = A$, giving
$\tilde\nu(\{\beta\}) = A/(\beta-1) = \lambda$, with no mass
elsewhere.  Since $\delta_0 > \delta_\infty$ (nontrivial
intermittency) and $\ln r < 0$, $A < 0$ and $\beta - 1 < 0$, so
$\lambda = A/(\beta-1) > 0$.

Therefore $\nu = \lambda\delta_b$ with $b = (\ln\beta)/k < 0$ and
$\lambda > 0$.  The generator $\log W$ is compound Poisson with
deterministic jump size~$b$ and rate~$\lambda$: this is the
log-Poisson distribution.
\end{proof}

\begin{corollary}[Principal cascade classes]\label{cor:partition}
The log-ID cascade family is partitioned by A1:
\begin{center}\footnotesize
\begin{tabular}{@{}llll@{}}
\toprule
Class & L\'evy triplet & A1 & Determinacy \\
\midrule
Log-Poisson & $\sigma^2=0$, $\nu=\lambda\delta_b$ & Holds
  & Determinate \\
Log-normal & $\sigma^2>0$, $\nu=0$ &
  Fails ($\delta_\infty=-\infty$) & Indeterminate \\
Log-stable & $\sigma^2=0$, $\nu=$power-law &
  Fails & --- \\
General log-ID & any other & Fails & --- \\
\bottomrule
\end{tabular}
\end{center}
\end{corollary}


% ── Proposition 2 ───────────────────────────────────────────────────
\begin{proposition}[Determinacy Dichotomy]\label{prop:determinacy}
The two principal cascade multiplier laws are distinguished by moment
determinacy:

\textup{(a)} If A1 holds within a cascade (log-Poisson regime), then
\[
  \ln\mathbb{E}[W^p]
  = (\gamma\ln r)\,p + C\ln r\,(1 - \beta^{p/k}),
\]
and the second term is $o(p)$.
The Carleman sum diverges.  $W$ is moment-determinate.

\textup{(b)} If the exponents are quadratic, $\zeta_p = c_1 p +
c_2 p^2$ (\textup{\cite{K62}}/log-normal regime), then
$\mathbb{E}[W^p] = \exp(\mu p + \sigma^2 p^2/2)$.  The Carleman sum
converges.  $W$ is moment-indeterminate: multiple distinct
distributions share the same moments.
\end{proposition}

\begin{proof}
\textit{Part~\textup{(a)}.}  By Lemma~\ref{lem:exponent}, $\zeta_p =
\gamma p + C(1-\beta^{p/k})$, so $\ln\mathbb{E}[W^p] = \zeta_p \ln r
= (\gamma\ln r)\,p + C\ln r\,(1-\beta^{p/k})$.  Since
$\beta^{p/k} \to 0$, the second term converges to $C\ln r$,
hence is $o(p)$.  Moment determinacy then follows from
Lemma~\ref{lem:determinacy}.

\textit{Part~\textup{(b)}.}
For log-normal~$W$ with parameters $(\mu, \sigma^2)$:
\[
  \mathbb{E}[W^{2p}] = \exp(2\mu p + 2\sigma^2 p^2),
  \qquad
  \bigl(\mathbb{E}[W^{2p}]\bigr)^{1/(2p)} = \exp(\mu + \sigma^2 p)
  \to \infty.
\]
The Carleman sum $\sum_{p=1}^\infty \exp(-\mu - \sigma^2 p)$ converges
(geometric series with ratio $e^{-\sigma^2} < 1$ for $\sigma^2 > 0$).
Therefore the moments do not uniquely determine~$W$.
\end{proof}


% ── Theorem 3 ───────────────────────────────────────────────────────
\begin{theorem}[Stability]\label{thm:stability}
Let $\{W_n\}$ be an i.i.d.\ multiplicative cascade with
log-infinitely-divisible generator and nontrivial intermittency.  If
A1 holds to within $\varepsilon$---that is,
\[
  \bigl|\delta_{p+k} - (1-\beta)\delta_\infty - \beta\delta_p\bigr|
  < \varepsilon
  \qquad\text{for all } p \in k\mathbb{N}_0,
\]
then, after the change of variables $u = e^{kx}$ mapping the L\'evy
measure to~$(0,1)$, the positive measure $\eta = (1-u)\,d\tilde\nu$
satisfies
\[
  W_1\!\left(\frac{\eta}{\|\eta\|},\;\delta_\beta\right)
  \leq
  \left(\frac{(1+\beta)^2\,|\ln r|}{|A|}\right)^{\!1/2}
  \!\cdot\sqrt{\varepsilon},
\]
where $\|\eta\| = \eta((0,1))$ is the total mass and $A = (\delta_0 -
\delta_\infty)\ln r$.  In particular, the cascade multiplier
distribution converges to log-Poisson as $\varepsilon \to 0$.
\end{theorem}

\begin{proof}
The proof reuses the central construction of
Theorem~\ref{thm:classification}: the change of variables $u = e^{kx}$
that maps the L\'evy measure to the compact interval $[0,1]$.

\textit{Step~1 (Approximate moments on $[0,1]$).}  By the
classification proof, the A1 condition at $p = mk$ maps to moments of
the signed measure $\rho = (u-1)\tilde\nu$ on $(0,1)$:
\[
  \int_{(0,1)} u^m\,d\rho(u) = A\beta^m
  \qquad\text{(exact A1).}
\]
If A1 holds to within $\varepsilon$, the same derivation gives
\[
  \left|\int_{(0,1)} u^m\,d\rho(u) - A\beta^m\right|
  \leq |\ln r|\,\varepsilon
  \qquad\text{for all } m \geq 0.
\]

\textit{Step~2 (Variance bound).}  Define the positive measure $\eta
= (1-u)\,d\tilde\nu$ on $(0,1)$, so $d\eta = -d\rho$.  The moments
of $\eta$ satisfy $\int u^m\,d\eta = |A|\beta^m + O(\varepsilon)$.
Evaluate at the test function $f(u) = (u-\beta)^2$:
\begin{align*}
  \int(u-\beta)^2\,d\eta
  &= \int u^2\,d\eta - 2\beta\int u\,d\eta + \beta^2\int d\eta \\
  &= \bigl(|A|\beta^2 + O(\varepsilon)\bigr)
    - 2\beta\bigl(|A|\beta + O(\varepsilon)\bigr)
    + \beta^2\bigl(|A| + O(\varepsilon)\bigr) \\
  &= O(\varepsilon).
\end{align*}
Precisely:
\[
  0 \leq \int(u-\beta)^2\,d\eta
  \leq (1+\beta)^2\,|\ln r|\,\varepsilon.
\]

\textit{Step~3 (Concentration and Wasserstein bound).}  By
Chebyshev's inequality applied to the positive measure~$\eta$:
\[
  \eta\bigl(\{u : |u-\beta| > \delta\}\bigr)
  \leq \frac{(1+\beta)^2\,|\ln r|\,\varepsilon}{\delta^2}
  \qquad\text{for all } \delta > 0.
\]
The total mass $\int d\eta = |A| + O(\varepsilon)$.  By
Cauchy--Schwarz:
\[
  \int|u-\beta|\,d\eta
  \leq \Bigl(\int d\eta\Bigr)^{1/2}
       \Bigl(\int(u-\beta)^2\,d\eta\Bigr)^{1/2}.
\]
Therefore
\[
  W_1\!\left(\frac{\eta}{\int d\eta},\;\delta_\beta\right)
  \leq \left(\frac{(1+\beta)^2\,|\ln r|\,\varepsilon}{|A|}
       \right)^{\!1/2}
  = K\sqrt{\varepsilon},
\]
where $K = \bigl((1+\beta)^2|\ln r|/|A|\bigr)^{1/2}$.

\textit{Step~4 (Propagation to the multiplier distribution).}  It
remains to show that closeness of the L\'evy measure implies closeness
of the distribution of~$W$.  We work entirely with $\log W$, which is
a compound Poisson random variable.

Let $\nu_0 = \lambda\delta_b$ be the unperturbed (log-Poisson) L\'evy
measure and $\nu_\varepsilon$ be the perturbed L\'evy measure, with
total mass $\lambda_\varepsilon$.  By Step~2,
$|\lambda_\varepsilon - \lambda| \leq C_1\varepsilon$ for a constant
$C_1$ depending only on the unperturbed parameters.

\textit{Step~4a (Coordinate change).}  The map $\varphi\colon (0,1]
\to (-\infty,0]$, $\varphi(u) = (\ln u)/k$, is Lipschitz on
$[\beta/2, 1]$ with constant $L_\varphi = 2/(k\beta)$.  A tail mass
estimate via Markov's inequality gives
$F_\varepsilon^u\bigl((0,\beta/2)\bigr) \leq 2K\sqrt{\varepsilon}/
\beta$.  Combining:
\[
  W_1(F_\varepsilon, \delta_b) \leq K'\sqrt{\varepsilon},
\]
where $K'$ depends only on the unperturbed parameters.

\textit{Step~4b (Coupling).}  Let
$N_0 \sim \mathrm{Poisson}(\lambda)$ and
$M \sim \mathrm{Poisson}(\lambda_\varepsilon - \lambda)$
be independent, with $N_\varepsilon = N_0 + M$.  Then
\[
  X_\varepsilon - X_0
  = \sum_{i=1}^{N_0}(J_i - b)
  + \sum_{i=N_0+1}^{N_0+M} J_i.
\]
The shared-jump contribution is $\leq \lambda K'\sqrt{\varepsilon}$;
the excess-jump contribution is $\leq C_1(|b|+K')\varepsilon$.
Therefore $W_1(X_\varepsilon, X_0) \leq (\lambda K' + C_1(|b| +
K'))\sqrt{\varepsilon}$ for $\varepsilon \leq 1$.

\textit{Step~4c (Drift perturbation).}  The drift $a$ is determined
by $\mathbb{E}[W] = 1$.  The function $g(x) = e^x - 1$ is
1-Lipschitz on $(-\infty,0]$, giving $|a_\varepsilon - a_0| \leq
C_2\sqrt{\varepsilon}$.

\textit{Step~4d (Pushforward under the exponential map).}  Let
$(Y,Z)$ be an optimal coupling of $\log W_\varepsilon$ and
$\log W_0$ with
$\mathbb{E}[|Y-Z|] \leq C_3\sqrt{\varepsilon}$.
Fix $R = |\ln\varepsilon|$ and split
\begin{multline*}
  \mathbb{E}[|e^Y - e^Z|]
  = \mathbb{E}\bigl[|e^Y - e^Z|\,
    \mathds{1}_{\{|Y|\leq R,\,|Z|\leq R\}}\bigr] \\
  + \mathbb{E}\bigl[|e^Y - e^Z|\,
    \mathds{1}_{\{|Y|>R \text{ or } |Z|>R\}}\bigr].
\end{multline*}
For the bulk term, the mean value theorem gives $|e^Y - e^Z| \leq
e^{\max(Y,Z)}|Y-Z|$.  Cauchy--Schwarz on the bulk gives
\[
  \mathbb{E}\bigl[e^{\max(Y,Z)}|Y\!-\!Z|\,
  \mathds{1}_{\text{bulk}}\bigr]
  \leq
  \bigl(\mathbb{E}[e^{2\max(Y,Z)}
  \mathds{1}_{\text{bulk}}]\bigr)^{1/2}
  \bigl(\mathbb{E}[|Y\!-\!Z|^2]\bigr)^{1/2}.
\]
The first factor is bounded by $C_4 =
(4\,\mathbb{E}[W_0^2])^{1/2} < \infty$, since $\zeta_2$ is finite
and, for $\varepsilon$ sufficiently small,
$\mathbb{E}[W_\varepsilon^2] \leq 2\,\mathbb{E}[W_0^2]$ by
continuity of the compound Poisson MGF in the L\'evy measure.  From
the coupling in Step~4b, $\mathbb{E}[|Y-Z|^2] \leq
C_5\,\varepsilon$.  For the tail term, Chernoff bounds for Poisson
tails~\cite{BLM13} give $\mathbb{P}(|Z| > R) \leq e^{-cR}$ for a
constant $c > 0$; by Cauchy--Schwarz, the tail contribution is
$O(\varepsilon^c) = o(\sqrt{\varepsilon})$.  Combining:
\[
  W_1(W_\varepsilon, W_0)
  \leq \mathbb{E}[|e^Y - e^Z|]
  \leq C(\beta,\lambda,k,r)\,\sqrt{\varepsilon},
\]
where $C(\beta,\lambda,k,r)$ is an explicit constant depending only
on the unperturbed log-Poisson parameters.

The coupling in Step~4b follows the canonical construction for
comparing compound Poisson variables via independent thinning; see
Barbour, Holst, and Janson~\cite{BHJ92}, Theorem~10.A.
\end{proof}


\begin{remark}[Why the stability proof is elementary]
The change of variables $u = e^{kx}$ maps the L\'evy measure to the
compact interval $[0,1]$.  On a compact set: (a)~the Hausdorff moment
problem is always determinate (Weierstrass approximation, no Carleman
condition needed); (b)~stability follows from a second-moment test and
Chebyshev's inequality; (c)~the resulting rate is
$O(\sqrt{\varepsilon})$, meaning the log-Poisson class is an open set
in the space of cascade multiplier distributions.  The stability
theorem inherits its simplicity from the classification theorem.
\end{remark}


% ══════════════════════════════════════════════════════════════════════
\section{Corollaries}
\label{sec:corollaries}
% ══════════════════════════════════════════════════════════════════════

\begin{corollary}[Conservation constraint]\label{cor:conservation}
If there exists an index $k_0 > 0$ such that $\zeta_{k_0} = z_0$ for
a known constant $z_0$ fixed by an exact conservation law, then
\[
  \gamma = \frac{z_0 - C(1-\beta^{k_0/k})}{k_0}.
\]
This reduces the observable parameters from two to one.
\end{corollary}

\begin{corollary}[Codimension identification]\label{cor:codimension}
If the most singular structures have Hausdorff codimension
$C_{\mathrm{geom}}$ and $C = C_{\mathrm{geom}}$, then $\beta$
alone determines the full exponent curve, the multifractal spectrum,
and the cascade distribution.
\end{corollary}

\begin{corollary}[Spectrum width]\label{cor:width}
The width of the multifractal spectrum is
\[
  \Delta h = h_{\max} - h_{\min} = \frac{C}{k}\,|\ln\beta|,
\]
where $h_{\max} = \gamma + (C/k)|\ln\beta|$ (at $p = 0$, most
regular) and $h_{\min} = \gamma$ (at $p \to \infty$, most singular).
\end{corollary}


% ══════════════════════════════════════════════════════════════════════
\section{Concluding remarks}
\label{sec:discussion}
% ══════════════════════════════════════════════════════════════════════

The results of this paper show that the hierarchical symmetry~A1
carries considerably more force than might be expected from its
appearance as a simple linear recurrence.  Within the i.i.d.\
multiplicative cascade framework, A1 is equivalent to the log-Poisson
class: it is the unique axiom that selects a single distribution from
the full log-infinitely-divisible family, and it does so stably.

Several directions remain open.

\begin{enumerate}
\item \textbf{Beyond i.i.d.}  The i.i.d.\ assumption on the cascade
  multipliers is standard but restrictive.  It would be of interest to
  determine whether the classification extends to stationary ergodic
  or Markovian multipliers, where the L\'evy--Khintchine machinery is
  no longer directly available.

\item \textbf{Boundary cases.}  The present results assume nontrivial
  intermittency ($C > 0$) and $\beta \in (0,1)$.  The boundary cases
  $\beta \to 0$ (maximal intermittency) and $\beta \to 1$ (vanishing
  intermittency) deserve separate analysis.

\item \textbf{Determination of~$k$.}  The hierarchy step~$k$ is
  treated as a given parameter.  In applications, $k$~must be
  estimated from data; a data-driven procedure for selecting~$k$ and
  quantifying its uncertainty would be valuable.

\item \textbf{Quantitative stability in applications.}  The stability
  bound (Theorem~\ref{thm:stability}) provides an explicit constant
  $K(\beta, \lambda, k, r)$.  Computing this constant for specific
  physical systems (e.g., fully developed turbulence with $\beta =
  2/3$, $C = 2$) would yield concrete tolerances for the degree to
  which A1 can be violated while remaining close to log-Poisson.
\end{enumerate}

% ══════════════════════════════════════════════════════════════════════
\begin{thebibliography}{99}
% ══════════════════════════════════════════════════════════════════════

\bibitem{Akh65}
N.\,I.\ Akhiezer,
\textit{The Classical Moment Problem and Some Related Questions in
Analysis},
Oliver \& Boyd, Edinburgh, 1965.

\bibitem{BHJ92}
A.\,D.\ Barbour, L.\ Holst, and S.\ Janson,
\textit{Poisson Approximation},
Oxford University Press, Oxford, 1992.

\bibitem{BM02}
J.\ Barral and B.\,B.\ Mandelbrot,
Multiplicative products of cylindrical pulses,
\textit{Probab.\ Theory Related Fields}\ \textbf{124} (2002),
409--430.

\bibitem{BLM13}
S.\ Boucheron, G.\ Lugosi, and P.\ Massart,
\textit{Concentration Inequalities: A Nonasymptotic Theory of
Independence},
Oxford University Press, Oxford, 2013.

\bibitem{Dub94}
B.\ Dubrulle,
Intermittency in fully developed turbulence: log-Poisson statistics
and generalized scale covariance,
\textit{Phys.\ Rev.\ Lett.}\ \textbf{73} (1994), 959--962.

\bibitem{DG96}
B.\ Dubrulle and F.\ Graner,
Possible statistics of scale invariant systems and the observation of
intermittency in turbulence,
\textit{J.\ Phys.\ II France}\ \textbf{6} (1996), 817--824.

\bibitem{KP76}
J.-P.\ Kahane and J.\ Peyri\`ere,
Sur certaines martingales de Benoit Mandelbrot,
\textit{Adv.\ Math.}\ \textbf{22} (1976), 131--145.

\bibitem{K62}
A.\,N.\ Kolmogorov,
A refinement of previous hypotheses concerning the local structure of
turbulence in a viscous incompressible fluid at high Reynolds number,
\textit{J.\ Fluid Mech.}\ \textbf{13} (1962), 82--85.

\bibitem{Man74}
B.\,B.\ Mandelbrot,
Intermittent turbulence in self-similar cascades: divergence of high
moments and dimension of the carrier,
\textit{J.\ Fluid Mech.}\ \textbf{62} (1974), 331--358.

\bibitem{SL94}
Z.-S.\ She and E.\ L\'ev\^eque,
Universal scaling laws in fully developed turbulence,
\textit{Phys.\ Rev.\ Lett.}\ \textbf{72} (1994), 336--339.

\bibitem{SW95}
Z.-S.\ She and E.\,C.\ Waymire,
Quantized energy dissipation and log-Poisson statistics in fully
developed turbulence,
\textit{Phys.\ Rev.\ Lett.}\ \textbf{74} (1995), 262--265.

\bibitem{SZ09}
Z.-S.\ She and Z.-X.\ Zhang,
Universal hierarchical symmetry for turbulence and general multi-scale
fluctuation systems,
\textit{Acta Mech.\ Sinica}\ \textbf{25} (2009), 279--294.

\bibitem{ST43}
J.\,A.\ Shohat and J.\,D.\ Tamarkin,
\textit{The Problem of Moments},
American Mathematical Society, Providence, RI, 1943.

\end{thebibliography}

\end{document}
